% abnt-header.tex
% Política ABNT para ABNT Engine — NBR 14724:2011
% ============================================================
%
% Este arquivo implementa as escolhas de **política** ABNT para o
% ABNT Engine. O `template.tex` é o motor estrutural (Quarto +
% Pandoc + abntex2); este arquivo é a camada de política ABNT.
%
% Decisões documentadas:
%   - ADR 001: BibLaTeX-abnt (bibliografia)
%   - ADR 002: NBR 14724:2011 + legenda 12pt / fonte 10pt
%
% As alterações neste arquivo DEVEM ser acompanhadas de:
%   1. Atualização de ADR correspondente em docs/decisions/
%   2. Atualização da matriz normativa em docs/research/
%   3. Validação visual (Camada 4-6 da estratégia de testes)
%
% ============================================================
% PACOTES ADICIONAIS PARA POLÍTICA ABNT DE FIGURAS
% ============================================================
% subcaption para subfiguras com tamanho ABNT
\usepackage{subcaption}

% ============================================================
% MACROS PARA FONTE DA ILUSTRAÇÃO
% (NBR 14724:2011 §5.8 + NBR 14724:2024 §5.8)
% ============================================================
% \fonte{...} já existe nativamente em abntex2 v1.8.
% Aqui criamos aliases mais explícitos para alinhamento com a norma.

% Fonte: Elaborado pelo próprio autor.
% Uso: \fonteElaborado
\newcommand{\fonteElaborado}{%
  \fonte{Elaborado pelo próprio autor.}%
}

% Fonte: Adaptado de <origem>.
% Uso: \fonteAdaptado{LAMPORT (1994, p. 50)}
\newcommand{\fonteAdaptado}[1]{%
  \fonte{Adaptado de #1.}%
}

% Fonte: retirado de <origem>.
% Uso: \fonteRetirado{XYZ (ano)}
\newcommand{\fonteRetirado}[1]{%
  \fonte{Retirado de #1.}%
}

% Fonte: traduzido de <origem>.
% Uso: \fonteTraduzido{XYZ (ano)}
\newcommand{\fonteTraduzido}[1]{%
  \fonte{Traduzido de #1.}%
}

% ============================================================
% LISTA DE ILUSTRAÇÕES (NBR 14724:2011 §4.2.1.9)
% ============================================================
% \listoffigures existe nativamente (pacote tocbibind).
% Aqui adicionamos título ABNT-padrão.
\newcommand{\listadeilustracoes}{%
  \chapter*{LISTA DE ILUSTRAÇÕES}%
  \addcontentsline{toc}{chapter}{LISTA DE ILUSTRAÇÕES}%
  \listoffigures%
}

% ============================================================
% SUBCAPTION (ABNT)
% ============================================================
% Tamanho ABNT: caption de subfigura = 12pt (mesmo do corpo),
% label de subfigura = "Figura 1a" (sem ponto final após label).
\captionsetup[subfigure]{
  labelfont={normalsize},
  textfont={normalsize},
  labelformat=simple,
  labelsep=space,
  font=normalsize
}

% ============================================================
% FIGURA (NBR 14724:2011 §5.8)
% ============================================================
% \caption já é "Figura N -- Título" (textendash) em abntex2 v1.9.6.
% A fonte (declarada via \fonte) é menor (10pt) por padrão abntex2.
% Esta seção é mantida como espaço para override futuro (ex.:
% alterar separador para hífen simples, etc.).

% ============================================================
% LISTA DE TABELAS (NBR 14724:2011 §4.2.1.8)
% ============================================================
% Gera capítulo pré-textual com índice de tabelas.
% Uso: \listadetabelas
\newcommand{\listadetabelas}{%
  \chapter*{LISTA DE TABELAS}%
  \addcontentsline{toc}{chapter}{LISTA DE TABELAS}%
  \listoftables%
}

% ============================================================
% LISTA DE QUADROS (NBR 14724:2011 §4.2.1.10)
% ============================================================
% Pandoc-Crossref não tem ambiente nativo "quadro"; usuário deve
% criar índice manual. Este comando está reservado para uso futuro.
% Uso: \listadequadros
\newcommand{\listadequadros}{%
  \chapter*{LISTA DE QUADROS}%
  \addcontentsline{toc}{chapter}{LISTA DE QUADROS}%
  % Pandoc-Crossref não gera \listofquadros; reservado para extensão.
  \textit{Lista de quadros não gerada automaticamente. Use lista manual.}%
}

% ============================================================
% TABELA (NBR 14724:2011 §5.9 + IBGE)
% ============================================================
% Macro \IBGEtab nativa em abntex2 v1.8.
% Aqui configuramos tamanho ABNT:
%   - Caption = 12pt (mesmo do corpo), formato "Tabela N -- Título"
%   - Fonte da tabela = 10pt (menor)
%   - Conteúdo = 12pt (mesmo do corpo)
\captionsetup[table]{
  labelfont={normalsize},
  textfont={normalsize},
  labelformat=simple,
  labelsep=endash,
  font=normalsize,
  justification=centering,
  singlelinecheck=false
}

% Alias para \IBGEtab com fonte.
% Uso: \tabelaIBGE{Título}{Fonte}{Conteúdo tabular}
% Mantém compatibilidade com abntex2 + adiciona política ABNT.
\newcommand{\tabelaIBGE}[3]{%
  \begin{table}[htb]%
    \centering%
    \caption{#1}%
    \IBGEtabfontereduzida%
    #3%
    \fonte{#2}%
  \end{table}%
}

% ============================================================
% EQUAÇÃO (NBR 14724:2011 §5.10)
% ============================================================
% Equações têm numeração sequencial automática (1), (2), (3)...
% NBR 14724 não tem regra explícita sobre tamanho/estilo; abntex2
% usa o default do LaTeX. Apenas criando alias semântico.
%
% Uso: \begin{equation}\label{eq:label} ... \end{equation}
% Ou inline: $...$ (matemática inline)

% ============================================================
% LISTAGEM (CÓDIGO) — NBR não tem regra específica
% ============================================================
% Pandoc-Crossref gera \begin{figure} por default para blocos de
% código. Para usar ambiente "listing" (pacote listings), adicionar
% flag --listings na linha de comando do Pandoc.
%
% Aqui apenas garantimos que caption setup de listagem (se existir)
% segue o padrão ABNT.
\captionsetup[listing]{
  labelfont={normalsize},
  textfont={normalsize},
  labelformat=simple,
  labelsep=space,
  font=normalsize
}

% ============================================================
% QUADRO (NBR 14724:2011 §5.10)
% ============================================================
% Pandoc-Crossref não tem ambiente "quadro" built-in; usuários devem
% usar tabela com nome customizado ou criar ambiente custom.
% Esta seção é mantida como espaço para override futuro.

